\section{Balanced harmonic remainders and their order derivatives}
\label{clc:sec:harmonic}

The scale alignment also identifies the correct unequal truncation
lengths. For $N\ge0$ put
\begin{equation}\label{clc:eq:R}
 \mathcal R_{N,s}(y)=\Li_s(-y)-\sum_{k=1}^{N}\frac{(-y)^k}{k^s}.
\end{equation}
The empty sum is zero. For $a>0$ and $\sigma\in\{1,-1\}$ write
\begin{equation}\label{clc:eq:Hcolored}
 H_{N,\sigma}^{(v)}(a)=\sum_{k=0}^{N-1}\frac{\sigma^k}{(k+a)^v}.
\end{equation}
This finite function is entire in $v$. At $a=1$, $\sigma=1$, it is the
usual generalized harmonic number $H_N^{(v)}$.

\begin{theorem}[Unequal-scale harmonic-remainder product]
\label{clc:thm:remainder}
For all integers $m\ge1$, $N\ge0$, and all complex $s,t$, with $W=s+t$,
\begin{align}
 &\int_0^\infty\mathcal R_{mN,s}(y)
            \mathcal R_{N,t}(y^{-m})\frac{\dd y}{y}
 = I_m(s,t;0)\notag\\
 &\quad-m^{-s-1}\biggl[
 \pi\sum_{r=1}^{m-1}c_r^{(m)}
 H_{N,\sigma_m}^{(W)}\left(\frac r m\right)
 +\sigma_m W H_{N,\sigma_m}^{(W+1)}(1)\biggr].
 \label{clc:eq:remainder}
\end{align}
The integral is absolutely convergent and entire in its two orders.
Every independent order derivative can be passed through it.
\end{theorem}
\begin{proof}
The elementary shift of the Lerch series, first for $0<y<1$ and then
by continuation along $y>0$, gives
\begin{equation}\label{clc:eq:Rshift}
 F_{N+1,s}(y)=-(-y)^{-N}\mathcal R_{N,s}(y).
\end{equation}
Take $A=mN$ in \eqref{clc:eq:I-def}. The two factors in
\eqref{clc:eq:Rshift} have powers $y^{-mN}$ and $y^{mN}$, so they
cancel; their signs give
\[
 \int_0^\infty\mathcal R_{mN,s}(y)
       \mathcal R_{N,t}(y^{-m})\frac{\dd y}{y}
 =\sigma_m^N I_m(s,t;mN).
\]
This also proves convergence and all order-holomorphic assertions by
Theorem~\ref{clc:thm:pair} and the weighted lemma.

The finite shift formula is
\[
 \Phi(\sigma,v,a+N)=\sigma^N
             [\Phi(\sigma,v,a)-H_{N,\sigma}^{(v)}(a)]
 \quad(\sigma=\pm1).
\]
It follows directly by removing the first $N$ terms; at $\sigma=1$
it is continued through the usual Hurwitz pole in the combined formula.
All arguments in \eqref{clc:eq:pair} increase by $N$ when $A=mN$.
The resulting factor $\sigma_m^N$ cancels the preceding sign, proving
\eqref{clc:eq:remainder}. All harmonic corrections are finite entire
functions of the orders, so no additional continuation issue remains.
\end{proof}

An explicit example is
\begin{align}
 &\int_0^\infty
 \left[-\log(1+y)+y-\frac{y^2}{2}\right]
 \left[-\log(1+y^{-2})+y^{-2}\right]\frac{\dd y}{y}\notag\\
 &\hspace{25mm}=\pi G-\frac38\zeta(3)-\pi+\frac12.
 \label{clc:eq:Rexample}
\end{align}
This is the case $m=2$, $N=1$, $s=t=1$. The two brackets must be kept
together. Expanding the product and integrating its monomials separately
would introduce divergent integrals and would not prove the identity.
The unequal truncation lengths $2$ and $1$ are dictated by the scales;
they are not interchangeable with two equal truncation lengths.

\subsection{Logarithmic harmonic derivatives}
For every $h\ge0$,
\begin{equation}\label{clc:eq:Hjets}
 \partial_v^hH_{N,\sigma}^{(v)}(a)
 =(-1)^h\sum_{k=0}^{N-1}
       \frac{\sigma^k\log^h(k+a)}{(k+a)^v}.
\end{equation}
Consequently the derivatives of \eqref{clc:eq:remainder} give explicit
identities for products of polylogarithmic order derivatives and finite
harmonic-logarithmic corrections, with no new infinite coefficients.
The prefactor $m^{-s-1}$ contributes powers of $-\log m$ and must be
differentiated as well.

For example, the first $s$ derivative of the correction bracket in
\eqref{clc:eq:remainder}, evaluated at any pair $(s,t)$, is
\begin{align*}
 m^{-s-1}\biggl\{&\pi\sum_r c_r^{(m)}
       [\partial_W H_{N,\sigma_m}^{(W)}(r/m)
                     -\log m\,H_{N,\sigma_m}^{(W)}(r/m)]\\
 &+\sigma_m[(1-W\log m)H_{N,\sigma_m}^{(W+1)}(1)
                  +W\partial_W H_{N,\sigma_m}^{(W+1)}(1)]\biggr\}.
\end{align*}
The $t$ derivative has no derivative of the prefactor. Thus the explicit
harmonic identities retain the same scale asymmetry as the raw
polylogarithm integral. Arbitrarily high derivatives follow from the
finite Leibniz rule rather than an unproved extrapolation.
